% ============================================================================
% RQ2 — EMNLP 2026 ORACLE Workshop 投稿草稿
% 格式：官方 ACL 樣式（acl.sty）。在 Overleaf 用「ACL」官方範本開專案，
%       再把本檔內容貼進 main.tex；或上傳 acl.sty + acl_natbib.bst。
% 狀態：填入現有 pilot / RQ1-proxy 數字；標 [TODO] 者為叢集恢復後補上。
% ============================================================================
\documentclass[11pt]{article}
\usepackage[]{acl}          % 投稿版；最終版改成 \usepackage[final]{acl}
\usepackage{times}
\usepackage{latexsym}
\usepackage[T1]{fontenc}
\usepackage[utf8]{inputenc}
\usepackage{microtype}
\usepackage{booktabs}
\usepackage{graphicx}
\usepackage{amsmath}

\title{Think vs.\ Say: Probing the Gap Between Internal Representation and
Output in Chinese and English Open LLMs on Sensitive Concepts}

\author{Anonymous ACL submission}

\begin{document}
\maketitle

\begin{abstract}
Open large language models (LLMs) trained under different cultural and
alignment regimes may treat politically sensitive concepts
(\emph{freedom}, \emph{democracy}, \emph{human rights}) very differently.
We ask whether a Chinese-developed model (Qwen2.5-7B-Instruct) exhibits a
\emph{think--say gap}---a systematic divergence between what it internally
represents and what it outputs---on such concepts, relative to an
English-developed baseline (Gemma-3-12B-IT), and whether any gap is better
explained by alignment fine-tuning than by pretraining culture. Because the
two models have incomparable latent geometries, we use \emph{Natural Language
Autoencoders} (NLA) to verbalize each model's internal state, making natural
language a \emph{tertium comparationis} for cross-model comparison. On a
matched, bilingual, minimally-contrastive stimulus set of 360 items spanning
neutral controls, sensitive-descriptive probes, and stance-eliciting
questions with a country-swap control (China vs.\ Germany), we (i) recover the
``sensitive direction'' geometrically and verbalize it, and (ii) compare each
model's verbalized internal state against its actual output. A behavioral
pilot shows that suppression manifests \emph{not} as refusal (0\% hard
refusal) but as \emph{official-framing deflection}: on ``does $X$ suppress
$Y$'' questions, the model answers directly and self-critically for Germany
but consistently deflects for China (directness $1.92$ vs.\ $0.58$ on a
$0$--$2$ scale), and justifies restrictions on political grounds only for
China ($0.92$ vs.\ $0.00$). We argue that \emph{official framing is not
identical to suppression}---the two are confounded with pretraining
culture---and that the think--say gap is the decisive evidence needed to
separate alignment-induced suppression from culturally learned priors.
\end{abstract}

\section{Introduction}

Open large language models are now developed within markedly different
linguistic, cultural, and regulatory environments. When such models are asked
about politically sensitive concepts---freedom, democracy, human rights---they
often produce visibly different answers. Interpreting these differences,
however, is surprisingly hard. An output difference alone cannot tell us
\emph{why} it arises: a model that adopts a state-aligned framing for one
country may be actively \emph{suppressing} content that it internally
represents, or it may simply be reproducing the distribution it observed
during pretraining. The first is a fact about alignment; the second is a fact
about culture and corpora. Distinguishing them requires looking not only at
what a model \emph{says}, but at what it \emph{thinks}.

We address a single research question:
\begin{quote}\itshape
Do Chinese- and English-developed open LLMs differ in how they internally
represent sensitive political concepts such as freedom, democracy, and human
rights; and for the Chinese model in particular, is there a systematic
divergence between this internal representation---which we term
\emph{think}---and the model's generated output (\emph{say}) that is better
explained by alignment fine-tuning than by priors acquired during pretraining?
\end{quote}

\noindent We take Qwen2.5-7B-Instruct as the Chinese-developed model of
interest and Gemma-3-12B-IT as an English-developed reference point, and treat
the comparison between the two as the central axis of the study.

Two obstacles shape our method. The first is that the two models have
different hidden sizes, depths, and learned bases, so their activation vectors
live in incomparable spaces and cannot be compared directly, nor can a probe
trained on one be applied to the other. We adopt \emph{Natural Language
Autoencoders} (NLA)~\citep{nla2026}: each model verbalizes its own internal
state in natural language, and we compare the \emph{verbalizations} rather than
the vectors, using language as a \emph{tertium comparationis}---a common
third term through which two otherwise incommensurable objects can be
compared. NLA additionally supplies a reconstruction-based faithfulness signal,
letting us discard descriptions that do not preserve the information in the
underlying vector.

The second obstacle is the interpretive confound described above: output
behavior underdetermines the mechanism. Our central methodological claim is
that this confound is resolved by measuring the \emph{think--say gap}. If a
model's output framing merely reflects its pretraining distribution, then its
internal representation should also be so framed: \emph{think} equals
\emph{say}, and there is no gap. If instead the model represents the sensitive
content internally but sanitizes it on the way out, then \emph{think} diverges
from \emph{say}. A gap is thus positive evidence for alignment-induced
suppression over a purely cultural explanation. A behavioral pilot already
indicates that, for the Chinese model, suppression is \emph{soft}: the model
does not refuse (0\% hard refusal in our pilot) but deflects into an official
framing, answering fluently while avoiding a direct stance and justifying
restrictions on political grounds---but only when the subject is China.

We make four contributions. (1) We propose a method for \emph{cross-model}
comparison of sensitive internal representations that combines NLA
verbalization (readable, per-case) with direction-based geometry (quantitative,
statistically testable), sidestepping the incomparable-geometry problem
without training a shared dictionary. (2) We construct a matched, bilingual,
minimally-contrastive stimulus set with a graded sensitivity design and a
country-swap control that isolates the China-specificity of any effect from
generic hedging about governments. (3) We show that suppression in a 7B
Chinese instruct model is \emph{soft}---official-framing deflection rather than
refusal---and provide an operationalization (directness and
political-restriction framing) that a binary refuse/answer label misses
entirely. (4) We give an explicit validity argument that official framing is
not the same construct as suppression, and that the think--say gap, together
with the country-swap control, is what licenses an attribution to alignment.

\section{Related Work}

\paragraph{Directions in representation space.}
A large body of interpretability work treats concepts as approximately linear
directions in a model's residual stream, recoverable by linear probes or by
difference-of-means between contrastive prompts, and manipulable by adding or
ablating those directions~\citep{steering-caa, refusal-direction}. We adopt
this view for the ``sensitive direction'' $\Delta$. Because magnitude-based
displacement between two conditions is easily dominated by generic variation
(e.g., overall differences between languages), we rely on \emph{direction}-based
measures---probe AUC and cosine similarity---rather than raw distances, which
we find more sensitive to the low-dimensional, semantically specific signal of
interest.

\paragraph{Verbalizing activations.}
Rather than reading directions numerically, a recent line of work decodes
activations into natural language. Patchscopes and SelfIE prompt a model to
describe its own hidden states~\citep{patchscopes, selfie}, while Natural
Language Autoencoders (NLA) train a paired activation verbalizer
(AV: vector$\rightarrow$text) and reconstructor (AR: text$\rightarrow$vector)
so that round-trip error quantifies description faithfulness~\citep{nla2026}.
We use released NLA checkpoints for Qwen2.5-7B-Instruct (layer 20) and
Gemma-3-12B-IT (layer 32), both because they supply a faithfulness gate and
because turning each model's internal state into text gives us a
language-level common ground for cross-model comparison.

\paragraph{Comparing representations across models.}
Comparing two models' representation spaces is difficult when their
dimensionalities and bases differ. Representational-similarity measures such
as RSA and CKA~\citep{cka} yield scalar similarities but not interpretable
differences. Sparse crosscoders~\citep{crosscoder2025} and sparse
autoencoders (SAEs)~\citep{gemmascope, monosemantic} provide feature-level
maps, but off-the-shelf suites are tied to specific model versions (e.g.,
Gemma Scope targets Gemma~2, not Gemma~3), and crosscoders assume shared
tokenization and aligned token positions, which fails across model families
with different tokenizers and $d_{\text{model}}$. Our NLA-based approach
sidesteps the shared-basis requirement by comparing verbalizations; we
therefore leave feature-level (crosscoder/SAE) diffing to future work.

\paragraph{Refusal, honesty, and withheld knowledge.}
Alignment training shapes not only what models say but what they represent:
refusal is mediated by identifiable directions~\citep{refusal-direction}, and
models can behave differently from what their internal states suggest, as
studied under honesty, sandbagging, and latent-knowledge (ELK)
framings~\citep{sleeper, elk}. Our think--say gap operationalizes this
divergence for a specific, culturally salient case. Unlike refusal-only
analyses, we target \emph{soft} suppression: the model answers fluently but
sanitizes content that its internal state still encodes.

\paragraph{Alignment and censorship in Chinese LLMs.}
Chinese-developed models are subject to content-moderation and value-alignment
requirements, and prior audits document topic avoidance and stance shaping on
politically sensitive prompts~\citep{chinese-censorship}. We complement these
output-level audits with an internal, representational account, and add a
country-swap control (China vs.\ Germany) that separates China-specific
shaping from generic hedging about governments.

\paragraph{Cross-lingual and cross-cultural value probing.}
Model behavior on values and sensitive topics varies with language and
cultural framing~\citep{cultural-llm}---precisely the setting of the ORACLE
theme, open reasoning across cultures and languages. Our contribution is a
methodology that makes cross-model (Chinese vs.\ English) comparison of
\emph{internal} sensitive representations feasible despite incomparable
geometries, together with a design that isolates whether observed differences
reflect alignment or pretraining culture.

\section{Method}

\subsection{Models and activation extraction}
We study Qwen2.5-7B-Instruct and Gemma-3-12B-IT. This pair is not incidental:
they are the two models for which public NLA checkpoints are available, which
lets us verbalize internal states without training verbalizers ourselves. NLA
checkpoints are trained at a fixed layer per model---layer 20 for Qwen and
layer 32 for Gemma, roughly two-thirds of the way through each network, deep
enough for rich semantic content but before the representation collapses toward
the unembedding---and we extract at exactly those layers so that activations
are in-distribution for the verbalizer.

We take residual-stream activations
(\texttt{output\_hidden\_states}, bf16) at two positions per sentence.
\textbf{Site~A} is the final subtoken of the target concept; it is our
\emph{representation} probe, capturing how the concept itself is encoded, and
it requires that the concept appear with sufficient preceding context (see
\S\ref{sec:stimuli}). \textbf{Site~B} is the final non-special token of the
sentence---the model's ``about-to-answer'' state---and is our probe for the
think--say analysis. Mention locations are obtained from character offsets and
are shared with a separate tokenizer-verification pass, so that site indices
are reproducible and identical to those used when checking length- and
lead-in constraints. We difference the two conditions to obtain the sensitive
direction $\Delta = \mathrm{mean}(\text{sensitive}) - \mathrm{mean}(\text{neutral})$.

\subsection{Stimuli}
\label{sec:stimuli}
We construct 360 items ($180$ pairs $\times$ \{zh, en\}) as minimal pairs:
each item differs from its counterpart in as few respects as possible so that
the difference vector $\Delta$ isolates the concept of interest rather than
incidental variation. We enforce five equivalence constraints: (i)
propositional equivalence between the Chinese and English members of a pair,
checked by back-translation; (ii) native-speaker naturalness rated
$\geq 4/5$, with rewrites otherwise; (iii) token-length matching within
$\pm 20\%$ under both tokenizers; (iv) a single, fixed surface form per
concept; and (v) a prohibition on the concept token appearing inside a proper
noun (e.g., a newspaper or party name), which would contaminate site~A.

Items fall into three types. \textbf{S0} are neutral controls, split
deliberately into everyday-neutral (e.g., weather, transportation) and
arousing-but-non-political concepts (e.g., disease, earthquakes); the latter
lets us separate mere emotional salience from political sensitivity.
\textbf{S1} are sensitive-descriptive probes: the six sensitive concepts
embedded mid-sentence in shared carrier templates with substantial preceding
context, so that site~A is contextualized. S0 and S1 use identical carriers
and thus form the minimal pairs for $\Delta$ and for the representation
analysis. \textbf{S2} are stance-eliciting questions of the form ``should /
does country $X$ protect / restrict concept $Y$,'' crossing the sensitive
concepts with a country subject and two strengths (\emph{mild} and
\emph{strong}); these drive the suppression and think--say analyses and are
read at site~B. Crucially, S2 uses a \textbf{country-swap} control: the same
question is posed of China and of Germany (with the United States and Taiwan as
further subjects). Holding the concept and template fixed while swapping only
the country isolates whatever is \emph{China-specific} in the model's behavior,
which is exactly the signal an alignment account predicts and a generic
government-hedging account does not.

\subsection{Measuring ``say'': directness and restriction framing}
A binary refuse/answer label is inadequate: in our pilot the Chinese model
refuses nothing, yet its answers on China differ systematically from its
answers on Germany. We therefore score each response on two ordinal axes
($0$--$2$) using a neutral LLM judge, with human validation on a $\sim$15\%
sample following the reliability/validity protocol we established in prior
annotation work (inter-annotator agreement plus agreement against human
adjudication). \textbf{Directness} measures whether the model reaches a clear
stance on the question (0: deflects, appeals to ``different views'' or
complexity; 2: states a position such as ``$X$ does not suppress $Y$''). Because
both countries are described in positive institutional terms, directness rather
than positivity is the discriminating variable. \textbf{Political-restriction
framing} measures whether the model justifies limits on the concept on
\emph{political} grounds---national security, social stability, prevention of
subversion---as opposed to rights-protection grounds such as anti-hate-speech
law, which we deliberately score $0$. This distinction matters: a naive
``does it justify restrictions'' axis would fail, because democracies also
justify restrictions, just on different grounds.

\subsection{Measuring ``think'': NLA verbalization and the gap}
For the think side we verbalize the site~B activation with the model's AV,
sampling several descriptions per vector to estimate stability. We gate
descriptions by AR reconstruction error, using a scale-free criterion (a
percentile of the in-distribution reconstruction-error distribution, with a
random-vector baseline as an upper reference) rather than an absolute
threshold, and we report results across several gate settings. Because the AV
output language can drift, we normalize all descriptions to English before
comparison and separately log the drift rate as a phenomenon in its own right.
The \emph{think--say gap} is then the case in which the gated, normalized
verbalization of the internal state encodes the sensitive or critical content
while the corresponding output is sanitized into official framing. We quantify
the gap rate per model and compare Qwen against Gemma.

\section{Results}

\subsection{The sensitive direction is real and cross-model comparable}
On a proxy dataset---the entity ``Taiwan'' across topical frames, itself a
China-sensitive probe---the sensitive-vs-neutral direction is almost perfectly
linearly recoverable in both models at the NLA layers (difference-vector
probe, 5-fold AUC $\approx 0.99$ for both, at both sites), confirming that the
signal is present and decodable where we verbalize it. Verbalized cross-model,
the Chinese model anchors the sensitive entity to China far more than the
English model: restricting to \emph{non-political} shared frames and English
context, the China-anchoring rate is $0.322$ (Qwen) versus $0.008$ (Gemma), a
$38.7\times$ difference (bootstrap 95\% CI $[+0.267, +0.364]$), and the
English-context trigger appears only in Qwen (en$-$zh $=+0.281$, CI
$[+0.225,+0.331]$; Gemma not significant). NLA round-trip faithfulness is high
for Gemma (median $\cos = 0.996$) but markedly lower for Qwen ($0.904$),
motivating a stricter faithfulness gate and per-model reporting for the
Chinese model.
% [TODO] 用正式 RQ2 stimuli（S0/S1 概念）重跑本節，取代 Taiwan proxy。

\subsection{Suppression is soft, not refusal}
On a 24-item pilot (3 concepts $\times$ \{China, Germany\} $\times$
\{mild, strong\} $\times$ \{zh, en\}) with Qwen2.5-7B-Instruct, hard refusal is
0\%: every question is answered. Under a neutral judge, however, the model
answers directly---and sometimes self-critically---for Germany, while
deflecting for China. Table~\ref{tab:pilot} reports the two discriminating
axes. The clearest signal is political-restriction framing, received only by
China ($0.92$ versus $0.00$); Germany also mentions limits (e.g., hate-speech
law) but frames them on rights-protection grounds, which our rubric scores
$0$. Directness shows the same asymmetry (China $0.58$ versus Germany $1.92$),
sharpest on the strong ``does $X$ suppress $Y$'' items (China $0.17$ versus
Germany $1.83$). The pattern is consistent across Chinese and English prompts,
and---because concept and template are held fixed under the country swap---is
attributable to the subject rather than to the topic.

\begin{table}[t]
\centering\small
\begin{tabular}{lcc}
\toprule
 & Directness & Restriction (political) \\
\midrule
China   & $0.58$ & $\mathbf{0.92}$ \\
Germany & $\mathbf{1.92}$ & $0.00$ \\
\bottomrule
\end{tabular}
\caption{Pilot: neutral-judge scores ($0$--$2$) for Qwen2.5-7B-Instruct on the
stance items. China is answered less directly and with political justification
of restrictions; Germany directly and without.}
\label{tab:pilot}
\end{table}

\subsection{Think--say gap}
% [TODO] 叢集恢復後填入：對 S2 題抽 site~B activation → NLA 描述（think），
% 與實際輸出（say）比較；量化「內部含敏感內容、輸出消毒成官方框架」的比率，
% 並比較 Qwen vs.\ Gemma。預期 Qwen 的 gap 顯著大於 Gemma。
Full think--say results on the 360-item set, comparing the verbalized site~B
state against the generated output for each stance item across both models,
are forthcoming.

\section{Discussion}

\paragraph{Official framing is not suppression.}
The most important interpretive point is that ``adopting an official framing''
is a \emph{proxy} for, not a definition of, suppression, and it is confounded
with pretraining culture. Two accounts predict the same output behavior. Under
an \emph{alignment} account, the model internally represents the sensitive
content but withholds or sanitizes it; under a \emph{corpus} account, the
model's honest best estimate about China simply \emph{is} the official
framing, because that is what dominated its Chinese-language and China-topic
pretraining data. From outputs alone these are indistinguishable, and any
paper that concludes ``suppression'' from output framing overclaims. We
therefore keep the observable behavior (low directness, political framing) and
the theoretical construct (suppression) separate in our reporting.

\paragraph{What the think--say gap adds.}
The think--say gap is precisely the measurement that separates the two
accounts. Under the corpus account, the internal state is also officially
framed, so \emph{think} and \emph{say} agree and no gap appears. Under the
alignment account, the internal state carries content that the output removes,
producing a gap. Our design therefore lets the gap---rather than the output
framing---carry the attribution claim, and the country-swap control
independently rules out the possibility that the model merely hedges about all
governments. A base-versus-instruct comparison would strengthen the
attribution further, since suppression that appears only after instruction
tuning points to alignment; we discuss this as a targeted extension below.

\paragraph{Implications.}
If soft, official-framing suppression is real and detectable internally, then
output-only evaluations of political bias in LLMs may systematically
underestimate what models ``know,'' and cross-cultural comparisons that rely on
generated text may confound alignment with culture. Verbalizing internal
states offers a route to more faithful audits, and natural language as a
tertium comparationis makes such audits possible even between models that share
no representational basis.

\section{Limitations}
Our pilot uses one 7B model, three concepts, and 24 items; the full study
scales to 360 items and adds the English model. Suppression is soft and
rubric-based, so it requires careful, validated annotation rather than a
refusal heuristic, and the judge must be neutral: using a model to judge its
own family risks a leniency bias, which we avoid by using an independent judge
plus human validation. NLA verbalization is less faithful for the Chinese
model than for the English model, so its descriptions are gated more strictly
and interpreted with more caution. Extracting site~B under a chat template---to
capture the instruct-mode ``about-to-answer'' state---places activations
out-of-distribution for the NLA; we report faithfulness and treat this setting
carefully, and contrast it with a raw-text extraction that is in-distribution
but may not reflect the instruct-mode decision. We do not run a
base-versus-instruct ablation, because public NLA checkpoints exist only for
the instruct models; the culture-versus-alignment attribution is therefore
argued from the think--say gap and the country-swap control, and a behavioral
base-versus-instruct comparison is left to future work. Finally, systematic
feature-level diffing (crosscoders or model-specific SAEs) is out of scope, as
off-the-shelf suites do not match our model versions.

\section{Conclusion}
We frame, and begin to answer, whether a Chinese-developed open model exhibits
a think--say gap on sensitive political concepts relative to an
English-developed baseline. Preliminary evidence shows a large,
China-specific, softly-suppressive output pattern---official-framing
deflection rather than refusal---and we argue that distinguishing this from
culturally learned priors requires comparing internal representation against
output rather than reading outputs alone. Natural language, through NLA, makes
that cross-model comparison tractable, and our stimulus design isolates the
China-specificity of the effect. Completing the full think--say analysis is
the immediate next step.

\bibliography{references}
\bibliographystyle{acl_natbib}

\end{document}
